\section{Odd cyclic Clifford constraints for the non-self-dual residual}
\label{sec:r155-odd-clifford}

The anticommuting stabilizers in
Proposition~\ref{prop:stabilizer-rotation} constrain every odd pivotal
rotation, including rotations with an arbitrary central target. The
result gives a uniform necessary condition for the residual ring
$\RR{4}{155}$. Its complete tested range passes this condition.

\begin{theorem}[Odd cyclic Clifford factorization]
\label{thm:odd-cyclic-clifford}
Assume the hypotheses of Proposition~\ref{prop:stabilizer-rotation} for
an object $Y$. Let $n\ge3$ be odd, let $Z$ be a simple of
$Z(\mathcal C)$, and let $E$ be the generalized pivotal rotation on
\[
 V=\operatorname{Hom}_{\mathcal C}(F(Z),Y^{\otimes n}),
 \qquad E^n=\theta_Z^{-1}I.
\]
Choose an $n$th root of $\theta_Z$ and use it to normalize $E$ to
$E_0^n=I$. Put $r_n=2^{(n-1)/2}$ and define integers $c_j$ by
\[
 c_j=\#\left\{(\epsilon_1,\ldots,\epsilon_{(n-1)/2})\in\{\pm1\}^{(n-1)/2}:
 \frac12\sum_{a=1}^{(n-1)/2}a\epsilon_a\equiv j\pmod n\right\},
\]
where division by two takes place in $\ZZ/n\ZZ$. If $m_j$ is the
multiplicity of $\zeta_n^j$ as an eigenvalue of $E_0$, then there are
nonnegative integers $b_j$ such that
\begin{equation}\label{eq:odd-clifford-convolution}
 m_j=\sum_{k\in\ZZ/n\ZZ}c_{j-k}b_k,
 \qquad \sum_jb_j=\frac{\dim V}{r_n}.
\end{equation}
The circulant convolution matrix $(c_{j-k})$ is invertible over $\QQ$,
so the integers $b_j$, when they exist, are uniquely determined.
\end{theorem}
\begin{proof}
Use the normalized left and right stabilizer actions and pair-insertion
endomorphism $a$ from Proposition~\ref{prop:stabilizer-rotation}.
Insertions at adjacent positions anticommute, while insertions on disjoint
pairs commute, and every insertion squares to the identity. Naturality
of generalized rotation in the $n-1$ unmoved tensor factors transports
one ordinary adjacent insertion to the next. The rotation definitions
and naturality are those of \cite{NSpiv} and Ng--Schauenburg,
\emph{Congruence subgroups and generalized Frobenius--Schur indicators},
\href{https://arxiv.org/abs/0806.2493}{arXiv:0806.2493}, Lemmas~2.5--2.7.
Define the cyclic boundary insertion by conjugation with $E$. Since
$E^n$ is scalar, conjugation has exact period dividing $n$.
Thus the $n$ operators $A_j$, indexed modulo $n$, satisfy
\[
 A_j^2=I,\qquad
 A_jA_k=-A_kA_j\quad(j-k=\pm1),\qquad
 A_jA_k=A_kA_j\quad(j-k\ne0,\pm1).
\]
For $n=3$, the three distinct pairs all anticommute. For $n\ge5$, any
pair can be rotated away from the cyclic boundary, reducing these
identities to the adjacent or disjoint insertion calculations.

The abstract algebra with these relations has the ordered monomials
$A_0^{e_0}\cdots A_{n-1}^{e_{n-1}}$, $e_j\in\{0,1\}$, as a basis.
Its commutator form over $\mathbb F_2$ is the adjacency matrix of an
$n$-cycle. Its radical is the line of the all-ones vector: the radical
equations are $e_{j-1}=e_{j+1}$, and $n$ is odd. Hence the algebra is
$M_{r_n}(\CC)\oplus M_{r_n}(\CC)$.

For a concrete model, use Clifford generators $\gamma_j$ and put
$A_j=i\gamma_j\gamma_{j+1}$, with cyclic indices. Replacing every
$A_j$ by $-A_j$ gives the other simple block. The cyclic permutation
of the $\gamma_j$ is an orthogonal transformation of determinant one.
Its spin lift implements $A_j\mapsto A_{j+1}$ on either block.
Choose the unique sign of that lift for which its $n$th power is one,
and call the resulting operator $U$. The nontrivial orthogonal
rotation angles are $2\pi a/n$, $1\le a\le(n-1)/2$. The exterior-algebra
model of the spin representation therefore gives eigenvalues
\[
 \zeta_n^{\frac12\sum_a a\epsilon_a},\qquad \epsilon_a\in\{\pm1\},
\]
with the exponent interpreted modulo $n$. This includes the sign
normalization $U^n=I$, and gives exactly the multiplicities $c_j$.
Both simple algebra blocks have this same spectrum.

The operators $E_0$ and $U$ implement the same algebra automorphism.
Their quotient $U^{-1}E_0$ consequently acts on the multiplicity spaces,
commutes with $U$, and has $n$th power one. Tensoring its cyclic
characters with the spin character proves
\eqref{eq:odd-clifford-convolution}. A different chosen $n$th root of
$\theta_Z$ shifts all cyclic labels, and reversing the rotation reverses
them; neither change affects feasibility.

The eigenvalues of the circulant matrix are the traces of the powers
of $U$. Every power of an odd cycle is a permutation with only odd
cycles, so it has no orthogonal eigenvalue $-1$. The spin-character
product of half-angle cosines is therefore nonzero. This proves
invertibility.
\end{proof}

\begin{proposition}[All-order path divisibility for $R^\prime(m)$]
\label{prop:rprime-path-divisibility}
Under the same stabilizer hypotheses, every simple multiplicity in
$Y^{\otimes n}$ is divisible by $2^{\lfloor(n-1)/2\rfloor}$ for all
$n\ge1$. For the based rings $R^\prime(m)$ of
\eqref{eq:nonselfdual-neargroup}, the corresponding infinite family of
fusion-coefficient divisibilities holds if and only if $m$ is even.
\end{proposition}
\begin{proof}
The $n-1$ ordinary adjacent insertions generate the twisted group algebra
whose commutator form is the adjacency matrix of a path with $n-1$
vertices. This matrix has rank $2\lfloor(n-1)/2\rfloor$ over
$\mathbb F_2$. Every simple module consequently has dimension
$2^{\lfloor(n-1)/2\rfloor}$. Restricting the insertion action to each
simple multiplicity space proves the first assertion.

Write $x=[X]$, $y=[Y]$ in $R^\prime(m)$. Direct multiplication gives
\[
 x^3=m(1+h+g+hg)+2m^2x+(2+2m^2)y.
\]
Thus the divisibility at $n=3$ requires $m$ even. Conversely, suppose
$m$ is even. The coefficients of $x^3$ and $x^4$ are even, and the
fusion relations give
\[
 x^4=2mx^3+4mx+2(1+h),\qquad
 x^n=2mx^{n-1}+4mx^{n-3}+4x^{n-4}\quad(n\ge5).
\]
If the required divisibilities hold for smaller positive powers, the
three terms on the right have respective two-adic divisibilities at least
\[
 2+\left\lfloor\frac{n-2}{2}\right\rfloor,\qquad
 3+\left\lfloor\frac{n-4}{2}\right\rfloor,\qquad
 2+\left\lfloor\frac{n-5}{2}\right\rfloor.
\]
Each is at least $\lfloor(n-1)/2\rfloor$. The cases $n=1,2,3,4$
start the induction. Duality proves the same statement for $y$.
\end{proof}

\paragraph{Exact scope for $\RR{4}{155}$.}
The anticommuting-stabilizer hypothesis holds for $h$ and each of
$X,Y$. The generalized rotation spectra from
Section~\ref{sec:155-complete-modular} provide the input to
\eqref{eq:odd-clifford-convolution}. Exact rational inversion checks
both noninvertibles, all $52$ central targets, and every odd order
$3,5,\ldots,23$: all $1144$ blocks have nonnegative integral quotient
multiplicities. At order three the spin multiplicities are $(0,1,1)$,
so the quotient is explicitly $b_j=\dim(V)/2-m_j$.
The all-order path condition likewise holds because $m=4$ is even.
These assertions give necessary constraints and exact surviving data;
they supply neither associators nor a categorification or exclusion.

The source-bound checker is
\filename{evidence/rigidity_20261017/rank6/r155/verify_odd_clifford.py}.
It pins the full generalized-rotation certificate and source tensor by
SHA256, reconstructs each odd spin spectrum, solves every rational
convolution exactly, and rejects nonintegral or negative quotients.
Its negative controls run under both ordinary and optimized Python.
From the repository root, run
\begin{verbatim}
python evidence/rigidity_20261017/rank6/r155/verify_odd_clifford.py \
  --repo .
\end{verbatim}
The replay verifies these finite character assertions, not the categorical
arguments of Theorem~\ref{thm:odd-cyclic-clifford}.
